\section[Conformally twisted spectral triples for $C^*$-dynamical systems]{Conformally twisted spectral triples\\ for $\boldsymbol{C^*}$-dynamical systems}
\label{Sec:ConfTwistedTrSp}

In the following theorem, we use the selfadjoint operator $D_u$ to construct spectral triples for the natural actions of the algebra $A$ (with its left action on $\Hh$) and the algebra $A^\text{op}$ (acting on the right of $\Hh$).
\begin{Theorem}
\label{Thm:Main}
Let $G$ be a \emph{compact} Lie group of dimension $n$ acting ergodically on a~unital $C^*$-algebra $A$, then using the unique $G$-invariant trace $\varphi_0$ of Theorem~{\rm \ref{Thm:ErgodAct}}, we write $\Hh_0 := \GNS(A, \varphi_0)$.

For any fixed $h \in \Aa^1$ and any $u \in [0,1]$, the data $(A, \Hh_0 \otimes \bigwedge^\bullet \LieG^*, D_u)$ with grading $\gamma$ defines an even $n^+$-summable spectral triple, where
\begin{itemize}\itemsep=0pt
\item
the representation $\pi$ of $A$ on $\Hh = \Hh_0 \otimes \bigwedge^\bullet \LieG^*$ is given by restriction of the left multiplication~\eqref{Eqn:ProdOmega};
\item
the unbounded operator $D_u$ is the unique selfadjoint extension of
\begin{gather*}
D_{u} = K_{u} d K_{-u} + K_{-u} d^* K_{u},
\end{gather*}
defined on the core $\Core = \Aa^1 \otimes \bigwedge^\bullet \LieG^*$, the operator $K_u$ being defined by~\eqref{Eqn:DefKu};
\item
the grading operator $\gamma$ is defined on degree $k$ forms by
\begin{gather*}
%\label{Eqn:DefGamma}
\gamma( a \otimes v_1 \wedge \cdots \wedge v_k) = (-1)^k ( a \otimes v_1 \wedge \cdots \wedge v_k).
\end{gather*}
\end{itemize}
For any fixed $h \in \Aa^1$ and any $u \in [0,1]$, the data $(A^\text{op}, \Hh_0 \otimes \bigwedge^\bullet \LieG^*, D_u)$ with grading $\gamma$ defines an even $n^+$-summable \emph{twisted} spectral triple, with the automorphism $\beta$ on $A$ given by $\beta(a) = e^{h u} a e^{-h u}$ -- we use this $\beta$ to define an automorphism on $A^\text{op}$.
\end{Theorem}

\begin{Remark}
The morphism $\beta$ def\/ined above preserves the multiplication of $A^\text{op}$. It also satisf\/ies the relation \emph{unitarity condition} (see \cite[equation~(3.4)]{TrSpTypeIII}) that is $\beta\big( (a^\text{op})^* \big) = (\beta^{-1}(a^\text{op}))^*$.
\end{Remark}

\begin{proof}
It is clear from the def\/inition of $\pi$ that $A$ is represented on $\Hh$ by bounded operators. The existence and uniqueness of the selfadjoint extension of $D_u$ is proved in Proposition~\ref{Prop:Du}, while the compact resolvent and f\/inite summability properties are shown in Corollary~\ref{Cor:FiniteSumma}.

We now prove that the commutator of $D_u$ with $a \in \Aa^1$ is bounded. To this end, we use the notations of Lemma~\ref{Lem:Formed} to decompose the operator $D_u$. We call
\begin{itemize}\itemsep=0pt
\item
Part (0) is the ``bounded part'' of $D_u$, that is the terms
\begin{gather*}
 - K_u \bigg( \frac{1}{2} \sum_{k,\alpha,\beta} c^{i_k}_{\alpha \beta} \otimes B^{i_k}_{\alpha,\beta} \bigg) K_{-u}
\qquad
\text{and}
\qquad
- K_{-u} \bigg( \frac{1}{2} \sum_{k,\alpha,\beta} \overline{c^{i_k}_{\alpha, \beta}} \otimes (B^{i_k}_{\alpha, \beta})^* \bigg) K_{u}.
\end{gather*}
\item
Part (I) consists of the terms
\begin{gather*}
K_u \bigg( \sum_{j} \partial _j \otimes T_j \bigg) K_{-u}.
\end{gather*}
\item
Part (II) consists of the terms
\begin{gather*}
K_{-u} \bigg( \sum_{j} \partial _j \otimes T_j^* \bigg) K_{u}.
\end{gather*}
\end{itemize}
Part (0) commutes with the left multiplication by $a \in \Aa^1$, and thus it does not contribute to the commutator. We therefore only need to estimate Parts~(I) and~(II) of $D_u( a' \omegat)$ for $\omegat = a \otimes v_1 \wedge \cdots \wedge v_k$, that is
\begin{gather*}
\sum_j \partial_j\big( a' a e^{-(n/2 - k)h u} \big) e^{(n/2 - (k+1))h u} \otimes T_j( v_1 \wedge \cdots \wedge v_k) \\
\qquad\quad{}
+ \sum_j \partial_j\big( a' a e^{(n/2 - k)h u} \big) e^{-(n/2 - (k-1))h u} \otimes T_j^*( v_1 \wedge \cdots \wedge v_k)\\
\qquad{}
= \sum_j \big(\partial_j( a') a e^{-(n/2 - k)h u} + a' \partial_j\big(a e^{-(n/2 - k)h u}\big)\big) e^{(n/2 - (k+1))h u}
\otimes T_j( v_1 \wedge \cdots \wedge v_k) \\
\qquad\quad{}
+ \sum_j \big(\partial_j( a') a e^{(n/2 - k)h u} + a' \partial_j\big(a e^{(n/2 - k)h u}\big)\big) e^{-(n/2 - (k-1))h u} \otimes T_j^*( v_1 \wedge \cdots \wedge v_k).
\end{gather*}
It follows from these considerations that
\begin{gather*}
[D_u, a'] \omegat = \sum_j \partial_j(a') a e^{-h u} \otimes (T_j + T_j^*)( v_1 \wedge \cdots \wedge v_k),
\end{gather*}
which is clearly a bounded function of $\omegat$ for any $a' \in \Aa^1$. Moreover, such $a' \in \Aa^1$ sends the core~$\Core$ of our selfadjoint operator $D_u$ to itself and following \cite[Proposition~A.1, p.~293]{TrSpPiCr-Paterson}, this suf\/f\/ices to ensure that $a' \in \Aa^1$ sends the domain of~$D_u$ to itself. The algebra $\Aa$ of Def\/inition~\ref{Def:TrSp} thus contains $\Aa^1$ and is dense in the $C^*$-algebra $A$. This completes the proof that $(A, \Hh, D_u)$ is a~$n^+$-summable spectral triple.

\looseness=-1
It remains to study its parity: it is clear from the def\/inition that $\gamma$ sends the core $\Core$ to itself and thus it leaves the full domain of the selfadjoint operator $D_u$ stable. Clearly, $\gamma$ distinguishes only between $\Hh_\text{even} := A \otimes \bigwedge^\text{even} \LieG^*$ and $\Hh_\text{odd} := A \otimes \bigwedge^\text{odd} \LieG^*$ and $\pi(a)$ leaves both spaces invariant, while $D_u$ is an odd operator. This proves that $(A, \Hh, D_u)$ with $\gamma$ is an even spectral triple.

The parity paragraph above applies \textit{verbatim} to the spectral triple constructed from the right action of $A^\text{op}$. The summability property is also conserved. It remains to investigate the bounded twisted commutators. Notice f\/irst that if $a', h \in \Aa^1$ then both right multiplications by~$a'$ and by~$\beta(a')$ leave the core $\Core$ of $D_u$ invariant and therefore the domain of $D_u$ is also stable under these right multiplication.

Using the decomposition of $D_u$ into Parts (0), (I) and (II), it appears that Part~(0) commutes with the right action of~$A^\text{op}$ and therefore does not contribute to the commutator. We treat Parts~(I) and~(II) separately. Keeping only Part (I) in the expression $D_u( \omegat \cdot a')$ for $\omegat = a \otimes v_1 \wedge \cdots \wedge v_k$, we get
\begin{gather*}
\sum_j \partial_j\big(a a' e^{-(n/2 - k)h u} \big) e^{(n/2 - (k+1))h u} \otimes T_j( v_1 \wedge \cdots \wedge v_k) \\
\qquad{}
= \sum_j \big(\partial_j( a) a' e^{-(n/2 - k)h u} + a \partial_j(a') e^{-(n/2 - k)h u}\big) e^{(n/2 - (k+1))h u}
 \otimes T_j( v_1 \wedge \ldots \wedge v_k) \\
 \qquad\quad{}
+ \sum_{j} a a' \partial _j\big(e^{-(n/2 - k)h u}\big) e^{(n/2 - (k+1))h u} \otimes T_j(v_1 \wedge \cdots \wedge v_k)\\
\qquad{}
= \sum_{j} \big(\partial _j(a) a' e^{- h u} + a \partial _j(a') e^{-hu} + a a' \partial _j\big(e^{-(n/2 - k)h u}\big) e^{(n/2 - (k+1))h u} \big)\\
\qquad\quad{}
 \otimes T_j(v_1 \wedge \cdots \wedge v_k).
\end{gather*}
We compare this expression to $D_u(\omegat) \beta(a')$, i.e.,
\begin{gather*}
\sum_j \partial_j\big(a e^{-(n/2 - k)h u} \big) e^{(n/2 - (k+1))h u} e^{h u} a' e^{-hu} \otimes T_j( v_1 \wedge \cdots \wedge v_k) \\
\qquad{}= \sum_j \partial_j(a) a' e^{-h u} \otimes T_j( v_1 \wedge \cdots \wedge v_k) \\
\qquad\quad{} + \sum_{j} a \partial _j\big(e^{-(n/2 - k)h u} \big) e^{(n/2 - (k+1))h u} e^{h u} a' e^{-hu} \otimes T_j( v_1 \wedge \cdots \wedge v_k).
\end{gather*}
In these two sums, the only terms that could lead to an unbounded contribution are those containing~$\partial _j(a)$, but these two terms cancel. At this point, we must perform the same computation on Part~(II) to make sure that the automorphism~$\beta$ is also suitable for this case. A~very similar computation proves that this it is indeed the case~-- the key property is that $e^{-(n/2 - k)h u} e^{(n/2 - (k+1))h u} = e^{- hu} = e^{(n/2 - k)h u} e^{-(n/2 - (k-1))h u}$~-- and thus the operator (def\/ined \textit{a~priori} only on~$\Core$)
\begin{gather*}
%\label{Eqn:TwistedComm}
D_u \pi^\text{op}\big( (a')^\text{op} \big) - \pi^\text{op}( \beta(a')^\text{op}) D_u,
\end{gather*}
where $\pi^\text{op}\big( (a')^\text{op} \big) = \RightMult_{a'} \otimes \id_{\bigwedge^\bullet \LieG^*}$, extends to a bounded operator on~$\Hh$.
\end{proof}
